\section{NN-A003: a CUE theorem with the exterior phase retained}
\label{sec:NN-A003}

\paragraph{Inputs and status.}
This result uses no instance of A, M, E, or a GCT hypothesis.  Its inputs are
the Haar eigenangle density for $U(N)$ (Weyl integration), the elementary
determinantal integral identity proved below, and elementary Fourier analysis.
The Haar character orthogonality used in the final phase estimate also appears
in the supplied paper 12, Theorem \texttt{thm:first-deriv-mean}.  This is a
proved form within the specified CUE model, not a transfer theorem for zeta
zeros.  In particular the arithmetic phase is not discarded on the grounds
that it is bounded.

\begin{theorem}[NN-A003: marked CUE labels and their correlations]
Let $U_N$ be Haar distributed in $U(N)$, let
$0\leq\theta_1<\cdots<\theta_N<2\pi$ be its eigenangles, and choose an index
$I$ uniformly from $\{1,\ldots,N\}$, independently after sampling $U_N$.
Fix $0\leq r<1$ and $a\in\mathbb R$.  Define the continuous exterior phase
\[
 \phi_{N,r}(x)=\operatorname{Im}\sum_{j=1}^N
       \log(1-r e^{i(x-\theta_j)}),
 \qquad
 W_{N,i}=i-\frac{N\theta_i}{2\pi}-a-
                 \frac{\phi_{N,r}(\theta_i)}{\pi}.
\]
The logarithm is its power-series branch in $r$, so there is no branch
ambiguity.  Set $\eta_{N,i}=\operatorname{sgn}\sin(\pi W_{N,i})$; a convention
at zero has no effect on the limits.  For fixed $j\geq1$, let
\[
 D_{j,N}=\frac{N}{2\pi}(\theta_{I+j}-\theta_I),
 \qquad \theta_{i+N}=\theta_i+2\pi.
\]
Then $(W_{N,I}\bmod2,D_{j,N})$ converges in distribution to $(V,D_j)$,
where $V$ is uniform on $\mathbb R/2\mathbb Z$, independent of $D_j$, and
$D_j$ is the distance to the $j$th point to the right of zero in the reduced
Palm sine process.  The latter process has kernel
\[
 K^{0}(x,y)=\frac{\sin\pi(x-y)}{\pi(x-y)}-
      \frac{\sin\pi x}{\pi x}\frac{\sin\pi y}{\pi y}.
\]
Consequently
\[
 \lim_{N\to\infty}\mathbb P(\eta_{N,I}=1)=\frac12,
 \qquad
 \lim_{N\to\infty}\mathbb E[\eta_{N,I}\eta_{N,I+j}]
       =\mathbb E[1-2d_2(j-D_j)],                 \tag{A003.1}
\]
where $d_2(x)=\min_{m\in\mathbb Z}|x-2m|\in[0,1]$.
In particular the adjacent agreement probability tends to
$1-\mathbb E d_2(1-D_1)$.  All assertions remain valid for $r=r_N$ in a fixed
compact subinterval of $[0,1)$, provided a subsequence is used if a limiting
phase law itself is requested; the displayed label and gap limits are
independent of that phase law.
\end{theorem}

\begin{proof}
Normalize circle measure to $d\theta/(2\pi)$ and let $P_N$ be the orthogonal
projection onto the span of $1,e^{i\theta},\ldots,e^{i(N-1)\theta}$.
The squared Vandermonde density is the determinant of the kernel of $P_N$.
Expanding its two determinants and integrating each coordinate shows directly
that, for a bounded multiplier $u$,
\[
 \mathbb E\prod_i u(\theta_i)
       =\det(P_NuP_N|_{\operatorname{ran}P_N}).       \tag{A003.2}
\]
Indeed each term in the determinant expansion is a product of the matrix
entries $\int u(\theta)e^{i(k-l)\theta}d\theta/(2\pi)$, and the permutation
signs are precisely the determinant signs.  Conditioning one uniformly
marked point to be $x$ gives the same formula with the rank $N-1$ projection
\[
 Q_{N,x}=P_N-|v_x\rangle\langle v_x|,
 \qquad v_x(\theta)=P_N(\theta,x)/\sqrt N.
\]
This conditional kernel follows either by the Schur complement of the
determinant with the marked row and column fixed, or by expanding that row
and column in (A003.2).  The marked point itself is uniform on the circle.

For any finite-rank projection $Q$ and multiplier $u$ with $|u|=1$, put
$A=QuQ|_{\operatorname{ran}Q}$.  Then
\[
 |\det A|^2=\det\bigl(I_Q-Qu^*(1-Q)uQ\bigr)
       \leq \exp\bigl(-\|(1-Q)uQ\|_{\rm HS}^2\bigr),       \tag{A003.3}
\]
because the eigenvalues of the positive contraction subtracted from $I_Q$
are in $[0,1]$, and $1-y\leq e^{-y}$.  For $Q=P_N$ its exponent equals
\[
 F_N(u):=\|(1-P_N)uP_N\|_{\rm HS}^2
       =\sum_{n\in\mathbb Z}\min(N,|n|)|\widehat u(n)|^2.   \tag{A003.4}
\]
This follows by applying $u$ to each Fourier basis vector and counting how
many translates by $n$ leave $\{0,\ldots,N-1\}$.  Removing one vector from
the projection changes this leakage by at most one.  To see this write
$P=Q+R$, where $R$ has rank one.  Orthogonality gives
\[
 F(P)-F(Q)=-\|RuQ\|_{\rm HS}^2+
                 \|(1-P)uR\|_{\rm HS}^2\in[-1,1].       \tag{A003.5}
\]

Consider a nonzero Fourier mode $k\in\mathbb Z$ of $W\bmod2$.
When $k$ is even, $e^{i\pi kI}=1$.  The exterior phase at the marked
eigenangle and all cyclic gaps are invariant under simultaneous rotation of
all eigenangles.  Averaging the factor $e^{-ikNx/2}$ over that rotation
therefore makes the Fourier mode exactly zero, also after multiplication
by any bounded function of cyclic gaps.

When $k$ is odd, condition on the marked point $x\in(0,2\pi)$.  Its index
is $1+$ the number of remaining points in $(0,x)$.  Up to a scalar of
modulus one, the conditional Fourier mode is (A003.2) for $Q_{N,x}$ and
\[
 u_x(\theta)=(-1)^{\mathbf1_{(0,x)}(\theta)}
        \exp\{-ik\,\operatorname{Im}\log(1-r e^{i(x-\theta)})\}.
\]
This function is smooth except for two jumps, at $0$ and $x$, each of
absolute size two.  Piecewise integration by parts twice gives
\[
 \widehat u_x(n)=\frac{J_0+J_x e^{-inx}}{2\pi i n}+O_{x,k,r}(n^{-2}),
 \qquad |J_0|=|J_x|=2.
\]
For fixed $x\notin2\pi\mathbb Z$ the cross sum
$\sum_{n\leq N}e^{inx}/n$ is bounded.  Thus (A003.4), including positive
and negative frequencies, gives
\[
 F_N(u_x)\geq\frac4{\pi^2}\log N+O_{x,k,r}(1).
\]
Equations (A003.3)--(A003.5) imply that the conditional odd mode tends to
zero.  Its absolute value is at most one, so dominated convergence over
the marked angle proves marginal equidistribution.  Derivative bounds of
the smooth factor are uniform when $r$ stays bounded away from one.

The same estimate proves independence from local gaps, which does not
follow from marginal equidistribution alone.  Fix $s>0$ and let $J$ be the
arc of length $2\pi s/N$ immediately after $x$.  For $|z|\leq1$ replace
$u_x$ by $w=u_x z^{\mathbf1_J}$.  This is a contraction multiplier, and
\[
 \|(w-u_x)Q_{N,x}\|_{\rm HS}^2
      \leq 4\operatorname{Tr}(Q_{N,x}\mathbf1_JQ_{N,x})\leq4s.
\]
For a contraction $w$, $I_Q-(QwQ)^*(QwQ)$ is the sum of
$Q(1-|w|^2)Q$ and $Qw^*(1-Q)wQ$, both positive.  Hence its determinant
bound has exponent at least $\|(1-Q)wQ\|_{\rm HS}^2$.  By the triangle
inequality this is at least
\[
 \bigl(\|(1-Q)u_xQ\|_{\rm HS}-2\sqrt s\bigr)_+^2,
\]
which still tends to infinity.  It follows uniformly on every circle
$|z|=b<1$ that the mixed conditional generating function of the odd
Fourier mode and $z^{\#J}$ tends to zero.  Cauchy coefficient extraction
then gives the same assertion with $\mathbf1_{\#J=l}$ for every fixed
$l$.  Since $D_{j,N}>s$ means $\#J\leq j-1$, this proves vanishing of
all mixed Fourier modes against these gap events.  The even modes already
vanish by rotation.  The gaps are tight, since their mean is exactly $j$
by cyclic summation of all gaps.

After scaling by $N/(2\pi)$ and removing a harmless phase, the conditional
kernel tends uniformly on each compact interval to $K^0$ displayed in the
theorem.  Determinantal generating functions on such an interval converge:
their Fredholm series converge term by term, and Hadamard's bound for
positive kernels with diagonal at most one bounds the $m$th term by
$ (2s)^m/m!$. The limiting $D_j$ has no atoms: a determinantal process
with continuous intensity has probability zero of a point at any
prescribed location, and $D_j=s$ requires such a point at $s$.
The local count formulas and tightness therefore give convergence of
the scaled $j$th gaps to $D_j$ at every survival threshold. The mixed Fourier
identities just proved characterize the joint limit as the product of
Haar measure modulo two and this gap law.

It remains to compare the phase at neighboring eigenangles.  The convergent
Fourier expansion gives
\[
 \sup_x|\phi_{N,r}'(x)|\leq
          \sum_{m\geq1}r^m|\operatorname{Tr}(U_N^m)|.
\]
Direct Fourier integration of the CUE two-point determinant gives
$\mathbb E|\operatorname{Tr}(U_N^m)|^2=\min(m,N)$; the same identity
follows by substituting $e^{im\theta}$ into (A003.4).  Cauchy--Schwarz
therefore bounds the expectation of this supremum by
$\sum_{m\geq1}r^m\sqrt m<\infty$, uniformly in $N$.
Tightness of $D_{j,N}$ shows that
$\phi_{N,r}(\theta_{I+j})-\phi_{N,r}(\theta_I)$ tends to zero in
probability.  With unwrapped indices and angles the exact difference is
\[
 W_{N,I+j}-W_{N,I}=j-D_{j,N}
       -\frac{\phi_{N,r}(\theta_{I+j})-\phi_{N,r}(\theta_I)}{\pi}.
\]
Finally, for uniform $V\bmod2$, the two square-wave signs at $V$ and
$V+v$ agree on a set of normalized measure $1-d_2(v)$.  Their correlation
is $1-2d_2(v)$.  The discontinuity boundaries have zero limiting Haar
measure, so the joint convergence proves (A003.1).
\end{proof}

\paragraph{What this changes.}
The CUE half-proportion and correlation formulas are justified within an
explicit model containing the correlated exterior phase.  The earlier
nonperiodic expression $1-\mathbb E|1-D_1|$ for adjacent agreement must be
replaced by $1-\mathbb E d_2(1-D_1)$; gaps larger than two cannot be silently
omitted.  To apply this to the arithmetic labels still requires a joint
arithmetic limit theorem for the label phase and local zero gaps.
Montgomery's restricted pair-correlation theorem supplies neither the odd
label modes nor the nearest-neighbor law.  This finding supplies no
algorithm for a Boolean counting problem and no separation of complexity
classes.

\paragraph{Geometric and representation-theoretic dictionary.}
At finite $N$, (A003.2) is an exterior-algebra determinant of a compressed
operator, derived from Weyl integration on the compact real form of
$\mathrm{GL}_N$.  The rank-one Palm operation is removal of the evaluation
line, and the leakage in (A003.3) is a Hilbert--Schmidt norm.  These are
precise Hilbert-space and representation-theoretic constructions; they do
not assert a scheme over $\mathbb F_q$, a finite Weil object, or Frobenius
purity.  The limiting sine projection has infinite rank.  No Grothendieck
archive cote is claimed for this new analytic argument.

\paragraph{Exact scope of the classical arithmetic pair theorem.}
The primary locator is Montgomery, \emph{The pair correlation of zeros
of the zeta function} (1973), Section~1, Theorem and equation~(2),
printed p.181; the Fourier identity (3)--(4) and explicit support
restriction are on p.182. Under RH, for each fixed $\epsilon>0$,
\[
 F(\alpha,T)=(1+o(1))T^{-2|\alpha|}\log T+|\alpha|+o(1)
 \quad (|\alpha|\le1-\epsilon),
\]
with the stated uniformity. The full continuation beyond this range
is presented as conjecture (11)--(12), pp.183--184. This controls
band-limited pair observables, while a neighboring-gap event requires
an empty interval and thus higher count information. The arithmetic
derivative and exterior phase are additional marks. Neither of these
missing pieces is supplied by that restricted pair theorem. This is
existing background, not a repeated consequence from the prior catalogue.
